\chapter[Representation by WIS]{Representation by Weighted Incidence Structures}\label{ch:sec:representation}

We have introduced weighted incidence structures as the intrinsic
invariant of universally optimal encoders.  But do all WIS actually
arise from real encoders?  Given abstract weight functions $w,\beta$
and an incidence relation $I$, when does there exist a deterministic
encoder whose multiplicities factor exactly as $n_{mc}=\beta(c)/w(m)$?
This is the \emph{representation problem}, and its answer is a complete
combinatorial characterisation.

We can now state the forward direction of the classification: every
universally optimal deterministic encoder induces a weighted incidence
structure.
\begin{theorem}[Representation Theorem]\label{ch:thm:representation}
Let $E:\Mset\times\Kset\to\Cset$ be a deterministic Shannon encoder, with uniform keys and a full-support prior $\pi$. If $E$ is universally optimal relative to $\pi$, then its induced data form a weighted incidence structure. More precisely, for every active ciphertext $c$ there is a unique number $\beta(c)>0$ such that
\[
 n_{mc}=\frac{\beta(c)}{\pi(m)}
 \qquad\text{for every }m\in\cloak(c).
\]
\end{theorem}
\begin{proof}
Fix an active ciphertext $c$. By Corollary~\ref{ch:cor:uo-factorisation}, the quantity $\pi(m)n_{mc}$ is independent of $m\in\cloak(c)$. Define its common value to be $\beta(c)$. It is positive because $c$ is active and every $m\in\cloak(c)$ has positive prior and positive multiplicity. Hence $n_{mc}=\frac{\beta(c)}{\pi(m)}$ for every incidence. With $I=\{(m,c):n_{mc}>0\}$ and $w=\pi$, both conditions of Definition~\ref{ch:def:wis} hold. Uniqueness of $\beta(c)$ follows from the displayed identity at any incident message.
\end{proof}

In particular, the induced multiplicities satisfy the cycle identities
of Theorem~\ref{ch:thm:cycle-factorisation} on every connected component
of the incidence graph.  The Representation Theorem is stronger than
that graph-theoretic compatibility statement in one essential respect:
it identifies the message weights with the prescribed prior $\pi$.
Conversely, the cycle identities alone yield only an abstract WIS
factorisation and do not imply universal optimality for a fixed prior.
Together with the Characterization Theorem (Theorem~\ref{ch:thm:realizability}),
we obtain a complete classification: a deterministic encoder is
universally optimal relative to $\pi$ if and only if its multiplicity
matrix factorizes as a WIS whose weights satisfy the integrality and
balance conditions of Theorem~\ref{ch:thm:realizability} and where $w$ is
proportional to $\pi$.
\begin{corollary}\label{ch:cor:uo-realizability}
A weighted incidence structure $\mathfrak W$ is induced by a universally
optimal deterministic encoder with prior proportional to $w$ if and only if
\begin{enumerate}
\item it satisfies the two conditions of Theorem~\ref{ch:thm:realizability};
\item $w$ is proportional to the prior $\pi$.
\end{enumerate}
\end{corollary}
\begin{proof}
Combine Theorem~\ref{ch:thm:representation} (forward direction),
Theorem~\ref{ch:thm:realizability} (realizability of the WIS), and
Corollary~\ref{ch:cor:fixed-encoder} (compatibility with the prior).
\end{proof}

\begin{corollary}[Criterion at fixed encoder]\label{ch:cor:fixed-encoder}
Let $E$ be a deterministic encoder with multiplicity matrix $N$. Then $E$ is universally optimal relative to a full-support prior $\pi$ if and only if $N$ admits a WIS factorisation whose message weights are proportional to $\pi$.
\end{corollary}
\begin{proof}
The forward implication is Theorem~\ref{ch:thm:representation}. Conversely, if $n_{mc}=\beta(c)/w(m)$ and $w=\lambda\pi$ for some $\lambda>0$, then $\pi(m)n_{mc}=\beta(c)/\lambda$ is constant on every cloak. Corollary~\ref{ch:cor:uo-factorisation} applies.
\end{proof}

The converse representation problem is therefore closed.

The representation theorem and the gauge characterisation of this chapter are the foundation of Part~I. They will be used in Chapter~\ref{ch:sec:graphs} to construct the logarithmic linearisation, and every WIS mentioned in the rest of the book---including our running examples---is an instance of this representation.


\section{Uniqueness and Gauge Equivalence}\label{ch:sec:gauge}

The representation theorem gives a factorisation. We next show that its
only ambiguity is componentwise scaling.

\begin{lemma}[Edge ratios]\label{ch:lem:edge-ratios}
Suppose a multiplicity matrix satisfies
\(n_{mc}=\frac{\beta(c)}{w(m)}=\frac{\beta'(c)}{w'(m)}\) on every
incidence.  Then, for each edge \(mc\),
\(\frac{w'(m)}{w(m)}=\frac{\beta'(c)}{\beta(c)}.\)
\end{lemma}

\begin{proof}
Cross-multiplication of the two factorisations gives the asserted
equality. \(\)
\end{proof}

\paragraph{What this theorem proves.}
The WIS factorisation $n_{mc}=\beta(c)/w(m)$ is not unique: if $(w,\beta)$ works, then $(\lambda w,\lambda\beta)$ also works.  The question is whether this scaling ambiguity is the only one.  This theorem proves that it is: the WIS weights are unique up to a single global scaling factor---the gauge.  This is what makes the quotient geometry of Chapter~\ref{ch:sec:representation} well-defined.

\begin{theorem}[Uniqueness of the WIS weights]\label{ch:thm:uniqueness}
Let \(\mathfrak W\) and \(\mathfrak W'\) be two weighted incidence
structures with the same incidence relation \(I\) and the same induced
multiplicity ratios
\(n_{mc}=\beta(c)/w(m)=\beta'(c)/w'(m)\).  On every connected component
\(H\) of their incidence graph that contains at least one edge, there is
a constant \(\lambda_H>0\) such
that

\[
w'(m)=\lambda_H w(m),\qquad
\beta'(c)=\lambda_H\beta(c)
\]

for every vertex of \(H\).
\end{theorem}

\begin{proof}
By Lemma~\ref{ch:lem:edge-ratios}, the ratio \(w'/w\) at a message
endpoint equals the ratio \(\beta'/\beta\) at the adjacent ciphertext
endpoint.  Along a path in a connected component, successive edge
equalities propagate one common positive value.  Denote it by
\(\lambda_H\). \(\)
\end{proof}

\begin{corollary}[Gauge equivalence]\label{ch:cor:gauge}
On a connected incidence graph, two WIS factorisations of the same
multiplicity matrix differ by one global positive scaling.  On a
disconnected graph, the scaling is independent on distinct connected
components.
\end{corollary}

\begin{quote}
\textbf{Mathematical intuition.}
The gauge transformation $w\mapsto\lambda w,\;\beta\mapsto\lambda\beta$ leaves every incidence ratio $n_{mc}=\beta(c)/w(m)$ unchanged.
The only observable quantities are the ratios, not the weights themselves.
\end{quote}

Thus the intrinsic datum is a gauge class of weights, not a preferred
normalization.

\paragraph{Why the quotient?}
The gauge transformation $w\mapsto\lambda w,\;\beta\mapsto\lambda\beta$ leaves every observable quantity---the multiplicities $n_{mc}$---unchanged.  The weights themselves are not identifiable from the encoder.  Why not fix a gauge by requiring, say, $\sum_m w(m)=1$?  Because any such choice is arbitrary and obscures the fact that the theory depends only on ratios.  The quotient space $V_I/\operatorname{Im}(C)$ is the natural information space: two vectors that differ by a gauge transformation encode the same cryptographic structure.

\paragraph{Running example: browser fingerprinting.}
The browser fingerprinting example (Section~\ref{ch:sec:browser-fingerprinting-and-anonymity}) illustrates the power of the representation theorem.  With $|\Mset|\approx10^6$ distinct browsers and $|\Cset|$ limited by measurable attributes, the incidence structure is highly non-uniform.  For a browser $m$ with fingerprint $c$, the ratio $w(m)/\beta(c)=1/n_{mc}$ gives the per-fingerprint probability weight, and the gauge freedom $w\mapsto\lambda w,\;\beta\mapsto\lambda\beta$ corresponds to rescaling all browser identities simultaneously---the only ambiguity is a single multiplicative constant.

The choice \(w=\pi\) made in
Definition~\ref{ch:def:induced-wis} is a convenient representative
supplied by the probabilistic model.
